\section{Prior work and claim boundary}

\paragraph{Exact H\'enon orbit algebra.}
Exact low-period equations, orbital sums, discriminants, reversibility
classes, and Galois groups for the area-preserving H\'enon map have been
developed in detail by Endler and Gallas
\citep{endlergallas2002,endlergallas2006sums}.  Polynomial bridges can make
apparently different orbit equations birationally or algebraically
equivalent \citep{brison2018}.  These results prevent us from treating the
period-five sextic or its ordinary permutation group as a new cover.  Our
linear inverse subresultant makes the equivalence explicit; the contribution
lies in the action-image singularity and the stability data above it.

\paragraph{Generating actions and equal-action saddles.}
Periodic action principles for reversible symplectic maps are classical
\citep{kookmeiss1989}.  In the quantum-H\'enon setting, equal-action saddles
already enter Stokes geometry and virtual turning-point constructions
\citep{shudoikeda2008}.  We therefore do not claim the first H\'enon action
or the first equal-action H\'enon pair.  Our setting differs in being a
closed exact-period-five family over the parameter line, with an exact
collision polynomial and a return-stability square class.

\paragraph{Critical-value geometry.}
Discriminants of critical values and their monodromy belong to the
Lyashko--Looijenga framework \citep{looijenga1974}.  In that convention, the
locus where distinct Morse points share a critical value is a Maxwell
stratum; some real catastrophe-theory conventions reserve ``Maxwell set''
for equal global minima.  We use the neutral phrase \emph{equal-action
node} and make no minima claim.  The generic occurrence of a node or a
symmetric critical-value group is not new.  What is specific here is the
explicit H\'enon polynomial and its chronological Hill decoration.

\paragraph{Hill and Kummer structures.}
The equality between a discrete action Hessian and a monodromy determinant
is a specialization of Hill's formula \citep{bolotintreschev2010}.  Kummer
covers and their braid monodromy are likewise established
\citep{artal2014}.  We use these as structural boundaries: neither mechanism
is claimed as new, and an ordinary node does not automatically supply a
Picard--Lefschetz representation.  The theorem proved below is the coupled
specialization summarized in Table~\ref{tab:boundary}.

\begin{table}[t]
\centering
\small
\begin{tabular}{p{0.30\linewidth}p{0.25\linewidth}p{0.34\linewidth}}
\toprule
Object & Status before this work & Exact increment here \\
\midrule
Period-five marker cover & known & birationality rederived as a control \\
Discrete H\'enon action & known mechanism & exact action plane model \(W_5\) \\
Equal-critical-value locus & generic prior art & explicit coprime factor \(P_9\) \\
Plane node & generic prior art & two-point/tangent proof on \(P_9\) \\
Hill formula & known theorem & chronological H\'enon specialization \\
Kummer cover & known mechanism & nonsquare class \([h_1h_2]\) \\
\bottomrule
\end{tabular}
\caption{The novelty boundary.  The contribution is a coupled exact
specialization, not any generic mechanism in isolation.}
\label{tab:boundary}
\end{table}
